\documentclass[11pt]{article}


\setlength{\parindent}{0pt}
\usepackage{xltxtra}
\usepackage{hyperref}
\hypersetup{hidelinks}
\usepackage{url}
\urlstyle{tt}
\usepackage{xcolor}
\definecolor{CVBlue}{RGB}{23,110,191}
\usepackage{calc}
\usepackage{graphicx}
\usepackage{tikz}
\usetikzlibrary{calc}
\usepackage{fontspec}
\usepackage{xeCJK}
\usepackage{enumitem}
\CJKsetecglue{} %% 取消中文与数字之间的间隙


%% 主文档字体设置
\setmainfont[
    Path = fonts/Main/,
    Extension = .otf,
    BoldFont = texgyretermes-bold.otf, % 加粗字体
]{texgyretermes-regular.otf} % 正文字体

% 中文字体设置
\setCJKmainfont[
    Path = fonts/hansans/,
    Extension = .ttf,
    BoldFont = NotoSansSC-Bold.ttf, % 加粗字体
]{NotoSansSC-Regular.otf} % 正文字体


\usepackage{relsize}
\usepackage{xspace}

% 使用 fontawesome（部分图标）
\usepackage{fontawesome} 

% A4纸，上下左右边距
\usepackage[
    a4paper,
    left=1.2cm,
    right=1.2cm,
    top=1.5cm,
    bottom=1cm,
    nohead
]{geometry}

\renewcommand{\baselinestretch}{1.5} % 行间距设为1.5

\usepackage{titlesec}
\usepackage{enumitem}
\setlist{noitemsep} % 取消列表项间的额外间距
%\setlist{nosep} % 取消所有垂直间距
\setlist[itemize]{topsep=0.25em, leftmargin=*}
\setlist[enumerate]{topsep=0.25em, leftmargin=*}

% --- 用于控制【不同项目之间】的垂直距离 ---
\newlength{\interProjectSpacing}
\setlength{\interProjectSpacing}{0.9em} % <--- 在此调整项目之间的距离
\newcommand{\projectsep}{\vspace{\interProjectSpacing}}

% --- 用于控制【项目标题】与下方【项目描述】的距离 ---
\newlength{\intraProjectTitleSep}
\setlength{\intraProjectTitleSep}{0.4em} % <--- 在此调整标题和描述的距离
\newcommand{\titlebreak}{\\[\intraProjectTitleSep]}

% --- 用于控制【项目描述】与下方【要点列表】的距离 ---
\newlength{\intraProjectListTopSep}
\setlength{\intraProjectListTopSep}{0.2em} % <--- 在此调整描述和列表的距离

% =======================================================================


\titleformat{\section}         % 定制 \section 命令 
{\large\bfseries\raggedright} % 将 section 标题设置为大号、粗体且左对齐
{}{0em}                      % 可用于为所有 section 添加前缀（如“章节...”）
{}                           % 可用于在标题前插入代码
[{\color{CVBlue}\titlerule}]  % 在标题后插入一条横线
\titlespacing*{\section}{0cm}{*1.6}{*1.2}



\begin{document}
\pagenumbering{gobble}

%%%% 利用tikz来定位照片
\begin{tikzpicture}[remember picture, overlay] 
    \node[anchor = north east] at ($(current page.north east)+(-2cm,-1.2cm)$) {\includegraphics[height=3cm]{avatar.jpg}};
  \end{tikzpicture}%
  %%%% 利用tikz来定位学校Logo，这里只在第一页显示
  \begin{tikzpicture}[remember picture, overlay] 
    \node[anchor = north west] at ($(current page.north west)+(0.5cm,+1.0cm)$) {\includegraphics[height=6cm]{zju.png}};
  \end{tikzpicture}%
\centerline{\LARGE\bfseries{XXX}} 

\centerline{\normalsize{\faPhone\ 187-5821-8369 \quad \faEnvelopeO\ \href{mailto:3230100000@zju.edu.cn}{3240102877@zju.edu.cn}}} 

%\centerline{\normalsize{\faGithubSquare\ \href{https://github.com/maksymilan}{https://github.com/maksymilan} \quad \faRssSquare\ 
%\href{https://maksymilan.github.io/}{https://maksymilan.github.io}}} 
    
\section{\makebox[\widthof{\faGraduationCap}][c]{\color{CVBlue}\faGraduationCap}\ 教育背景}    
\textbf{浙江大学} \hfill 2024.09 -- 至今\\[0.5em] % 标题和正文间加一点距离
数学科学学院\quad 统计学\quad 大二\quad GPA 3.99/4.3
\begin{itemize}[nosep]
    \item 相关课程：《数学分析》、《高等代数与解析几何》、《概率论》、《数理统计》、《随机过程》、《统计学习》
\end{itemize}

\section{\makebox[\widthof{\faUsers}][c]{\color{CVBlue}\faUsers}\ 相关经历}

% --- 第一个项目 ---
% 将标题行末尾的 \\ 替换为 \titlebreak 命令

\textbf{复杂统计模型假设检验的多层准蒙特卡洛（MLQMC）框架功效评估} \hfill 2025.09 -- 2026.02 \titlebreak
项目描述：针对工业智能化中高维复杂模型假设检验的难题，项目旨在突破统计功效失真与计算效率低下的权衡瓶颈，提出了一种集高精度、高效率与强稳定性于一体的通用推断框架，相关成果发表于\quad \textbf{Expert Systems with Applications}（第一作者）。
% 在 itemize 的选项中，使用 topsep=\intraProjectListTopSep 来控制上边距
\begin{itemize}[nosep, topsep=\intraProjectListTopSep]
    \item \textbf{理论与算法创新}：提出融合 \textbf{随机化 Sobol 序列}、\textbf{分层敏感性分配} 与 \textbf{Logistic 平滑} 的多层准蒙特卡洛评估框架，解决了序贯蒙特卡洛（SMC）在 $O(10^3)$ 高维参数空间中的粒子退化与方差失控问题。
    \item \textbf{工程与性能优化}：成功抑制了小样本与非高斯噪声条件下的显著性水平偏移（FPR）；结合拟蒙特卡洛的高维均匀性，大幅缩减了尾部概率（极低 $\alpha$ 值）估算所需的采样量与计算耗时。
    \item \textbf{核心产出与应用}：验证了框架在半导体缺陷检测（第一类错误率 $< 0.1\%$）及金融风险监控中的实用性，实现了高维统计推断、方差缩减、平滑函数建模到工业级实时响应的验证。
\end{itemize}

% 使用 \projectsep 命令来分隔两个项目
\projectsep

\textbf{基于 JPL 历表与月球借力的绕日返回轨道设计与优化}（《数学软件》课程期末项目） \hfill 2026.06 -- 2026.07 \titlebreak
项目描述：基于钱学森《星际航行概论》拼接圆锥曲线理论，扩展并构建了结合月球引力弹弓效应（Swing-by）与 JPL Horizons 真实历表数据的工程级行星际轨道优化与四体（太阳-地球-月球-探测器）数值动力学仿真系统。
\begin{itemize}[nosep, topsep=\intraProjectListTopSep]
    \item \textbf{动力学建模与数值积分}：基于 \textbf{Python/SciPy} 实现了 Sun-Earth-Moon-Rocket 四体受力模型，采用 \textbf{Velocity-Verlet 辛积分器} 解决了高维非线性轨道演化的能量漂移问题，实现一年期演化能量相对漂移 $<10^{-6}$。
    \item \textbf{发射窗口优化与借力解析}：构建了月球作用域（SOI）内的双曲线偏折坐标变换算法；设计内外层双重网格扫描与优化算法，在 2026 年 365 天窗口内成功搜索出最小总速度增量 $$\Delta v_{\mathrm{total}}$$ 的最优发射日期与轨道参数。
    \item \textbf{工程对接与可重现性}：接入 \textbf{JPL Horizons API} 获取真实天体状态向量，完成数值轨道的真实历表残差校验（位置残差 $\le 6000\text{ km}$）；编写 \textbf{Makefile} 实现了涵盖数据抓取、仿真计算、分析绘图到轨道动画渲染的全自动流水线。
\end{itemize}

\section{\makebox[\widthof{\faCogs}][c]{\color{CVBlue}\faCogs}\ 技术栈}
\begin{itemize}[nosep]
    \item \textbf{编程语言：} Python, R
    \item \textbf{开发工具：} LaTex, Git
\end{itemize}
\section{\makebox[\widthof{\faGraduationCap}][c]{\color{CVBlue}\faList}\ 获奖情况}
\begin{itemize}
    \item 2024-2025学年浙江大学二等奖学金 \hfill 2025.10
    \item 2024-2025学年浙江大学学业优秀标兵 \hfill 2025.10
    \item 2024-2025学年浙江大学对外交流标兵 \hfill 2025.10
    \item 2026年美国大学生数学建模竞赛H奖(Honorable Mention) \hfill 2026.02
\end{itemize}
    
\section{\makebox[\widthof{\faInfo}][c]{\color{CVBlue}\faInfo}\ 其他}
\begin{itemize}[parsep=0.5ex]
    \item \textbf{英语水平：} TOFEL 110, CET-4 605, CET-6 602
    \item \textbf{交流经历：} 不列颠哥伦比亚大学暑期课程项目（数据科学方向） \hfill 2025.07-2025.08
\end{itemize}
\end{document}
